噪声不会因为你换了更好的编码就消失。每一次信道使用，符号都可能被打坏，这是物理事实。凭直觉，剩下的选择只有两个：认下一个正的错误率，或者不停堆冗余把它压低——而冗余堆到最后，速率会趋于零。

Shannon 的结论把这个二选一取消了。存在一个数 \(C\)，只要传输速率低于它，块错误概率就可以做到任意小；付出的代价仅是速率停在 \(C\) 以下，不是无限膨胀的冗余。噪声照旧每次都在，整块消息却可以几乎从不出错。

这条结论的另一半同样要紧：速率一旦越过 \(C\)，错误概率就压不下去了，而且在标准情形下会被推向 1。所以 \(C\) 不是一句鼓励，是一道墙。

\begin{center}
\begin{tikzpicture}[x=1cm,y=1cm,font=\small,>=stealth]
  \fill[blue!8] (0,0) rectangle (4.6,1.0);
  \fill[red!8] (4.6,0) rectangle (9.2,1.0);
  \draw[->,black!55] (0,0) -- (9.7,0) node[right,font=\footnotesize] {速率 \(R\)};
  \draw[black!70,thick] (4.6,-0.18) -- (4.6,1.35);
  \node[above,font=\footnotesize] at (4.6,1.35) {容量 \(C\)};
  \node[font=\footnotesize,align=center,black!75] at (2.3,0.5) {错误概率可压到任意小\\买单的是速率，不是无限冗余};
  \node[font=\footnotesize,align=center,black!75] at (6.9,0.5) {错误概率有正的下界\\继续改进编码没有用};
  \node[below,font=\footnotesize] at (0,-0.1) {\(0\)};
\end{tikzpicture}

\emph{整章讨论的就是这条竖线：它的位置由信道决定，两侧的行为质地完全不同。}
\end{center}

用一个数校准一下尺度。翻转概率 \(0.1\) 的二元对称信道，最朴素的三重重复加多数表决把单比特错误率压到约 \(2.8\%\)，速率从 1 降到 \(1/3\)。而这条信道的容量是 \(1-h_2(0.1)\approx0.531\) 比特每次使用——既不是重复码的 \(0.333\)，也不是未编码时的 \(1\)。两个数字之间的空隙，就是半个世纪的编码理论在填的东西。

\subsection{怎么读这一章}

第一次读，最值得花力气分清的是三个量。翻转概率描述一次传输有多吵；互信息描述某一种输入用法能传多少，它随你的输入分布变；容量则是在所有允许的用法里能拿到的最好成绩，只认信道本身。混淆前两个会以为容量随策略浮动，混淆后两个会把某次实测的互信息当成上限。

还有两处习惯性误读值得提前打预防针。编码定理谈的是块长趋于无穷时整块消息的错误概率，不承诺任何一位不出错，也不给出块长 \(n=1000\) 时的具体成绩；而读任何一个容量数值之前，都该先确认单位是比特每次使用、每实维度、每复符号还是每秒。

前一章\hyperref[ch:057]{``无损信源编码：熵为什么是压缩极限''}处理的是把数据变小，本章处理的是把数据送到；两条极限最后会在分离定理里碰头。

\subsection{本章篇目}

以下按概念依赖排列。若带着具体问题进来，也可以先打开最接近当前困惑的一篇，再顺着篇内链接回补。

\begin{itemize}
\tightlist
\item \hyperref[sec:07-Information-Theory/05-Channel-Coding/01]{``信道模型与一次传输的信息量''}：把信道、输入策略与联合分布拆开，说明为什么该量的是互信息而不是输出熵；
\item \hyperref[sec:07-Information-Theory/05-Channel-Coding/02]{``二元对称信道：从翻转概率到容量''}：算出 \(1-h_2(p)\)，并指出最坏的翻转概率是 \(0.5\) 而不是 \(1\)；
\item \hyperref[sec:07-Information-Theory/05-Channel-Coding/03]{``信道容量与有噪编码定理''}：一个单字母优化为什么能管长块的命运，可达与逆定理的骨架，以及分层设计所依赖的分离定理；
\item \hyperref[sec:07-Information-Theory/05-Channel-Coding/04]{``高斯信道：功率约束与连续容量''}：连续输入必须带着约束才有容量，以及 \(\frac12\log_2(1+\mathrm{SNR})\) 的两段脾气。
\end{itemize}

读完这四篇，回头看任何一个``提升吞吐''的工程方案，判断会多一个维度：它是在把系统推向那条竖线，还是在试图跨过去。前者值得投入，后者无论多聪明都不会成功。
